\chapter{Proofs for Approximate Quantitative Inference}
\label{annex:p15-approximate-quantitative-inference}

\begin{annexitem}{annex:p15.approximation.certificates}{thm:p15-approximation-certificates}{Intervals and comparison margins}
\begin{proof}
The first two implications follow by transitivity: \(f\geq L>c\) or
\(f\leq U<c\). When \(c\in[L,U]\), the possible value \(f=c\) satisfies
the interval premise and refutes either strict conclusion. Writing
\(f=\widehat f+e_f\), \(g=\widehat g+e_g\), the error in the difference
is \(e_f-e_g\), bounded by \(E_f+E_g\). The stated positive margin
leaves a positive exact difference. Finally
\(|f|\geq|\widehat f|-E_f>0\); division gives the relative bound.
\end{proof}
\end{annexitem}

\begin{annexitem}{annex:p15.approximation.log-bound}{thm:p15-approximation-log-bound}{Multiplicative rounding budgets}
\begin{proof}
Nearest integer selection gives \(|hj-\log_bx|\leq h/2\).
Exponentiating in base \(b>1\) yields the single-input inequalities,
including either tie choice. The base-\(b\) logarithm of the product
ratio is \(\sum_i a_i(\log_bQ_h(x_i)-\log_bx_i)\), whose absolute
value is at most \((h/2)\sum_i|a_i|\). Exponentiate again.
Negative powers are allowed because all inputs are positive.
The midpoint of consecutive representative logarithms corresponds to
\(b^{h(j+1/2)}\), their geometric mean. No independence was used.
\end{proof}
\end{annexitem}

\begin{annexitem}{annex:p15.approximation.lumping}{thm:p15-approximation-lumping}{The information discarded by a mean}
\begin{proof}
Weighted summation gives \(\sum_iw_i(a+bx_i)=a+b\mu\).
For a twice differentiable function, the fundamental theorem of calculus
applied twice along the segment from \(\mu\) to \(x_i\) gives
\[
 f(x_i)=f(\mu)+f'(\mu)(x_i-\mu)
 +(x_i-\mu)^2\int_0^1(1-t)f''(\mu+t(x_i-\mu))\,dt.
\]
Indeed, for \(F(t)=f(\mu+t(x_i-\mu))\), integrate
\(F'(t)=F'(0)+\int_0^tF''(s)\,ds\) from zero to one and exchange
the continuous integrals on the compact triangle. The weighted linear
terms vanish. Each remainder is at most \(M(x_i-\mu)^2/2\) in absolute
value, proving the bound. Expanding squares gives
\(\sum_iw_i(x_i-\mu)^2=\sum_iw_ix_i^2-\mu^2\).
Equal weights at zero and two give mean squared output two, whereas the
square of their mean is one. An exact input mean need not preserve a
nonlinear output. For a uniform selection \(J\) among \(N\) regions,
\(E[Nn_J]=\sum_{i=1}^Nn_i\); this proves the population example's
unbiasedness claim without bounding any particular realized error.
\end{proof}
\end{annexitem}

\begin{annexitem}{annex:p15.approximation.energy-average}{thm:p15-approximation-energy-average}{Momentum and kinetic energy}
\begin{proof}
Expansion gives
\[
 \sum_i m_i(v_i-\overline v)^2
 =\sum_i m_iv_i^2-2\overline v\sum_i m_iv_i+M_0\overline v^2
 =\sum_i m_iv_i^2-M_0\overline v^2.
\]
All terms on the left are nonnegative, and all vanish exactly when each
\(v_i=\overline v\), since every mass is positive. Multiply by one
half. The defining equation \(M_0\overline v=\sum_i m_iv_i\) proves
momentum preservation. For two unit masses with velocities six and four,
the difference is \(\tfrac12[(6-5)^2+(4-5)^2]=1\).
\end{proof}
\end{annexitem}

\begin{annexitem}{annex:p15.approximation.oscillators}{thm:p15-approximation-oscillators}{An exact period within the chosen force model}
\begin{proof}
Differentiation verifies the linear solution
\(x(t)=a\cos(\omega t)+b\sin(\omega t)\), \(\omega=\sqrt{k/m}\).
Every initial position and velocity chooses one pair \(a,b\). The
difference of two solutions with the same data has constant energy
\(m\dot x^2/2+kx^2/2\), initially zero, and is therefore zero.
A nonzero sinusoid has least positive period \(2\pi/\omega\).

For the nonlinear force, put \(V(x)=k|x|^{n+1}/(n+1)\).
Differentiating \(m\dot x^2/2+V(x)\) along a solution gives zero;
\(n>0\) makes the potential continuously differentiable even at zero.
At amplitude \(A\) the energy is \(kA^{n+1}/(n+1)\), so separation
during travel from zero to \(A\) gives quarter-period
\[
 \sqrt{\frac{m(n+1)}{2k}}\int_0^A
       \frac{dx}{\sqrt{A^{n+1}-x^{n+1}}}.
\]
Substitute \(x=Au\). The factor outside the resulting integral is
\(A^{(1-n)/2}\). Near one, the mean value theorem makes
\(1-u^{n+1}\) comparable to \(1-u\), whose inverse square root is
integrable; away from one the integrand is bounded. Reflection at the
turning points and through zero supplies four equal travel intervals.
This quadrature also constructs the motion: its inverse on an open travel
interval has the stated velocity, differentiation gives the force equation,
and the reflected position, velocity and acceleration agree at the joins.
The motion crosses zero with nonzero speed and reaches each turning point
with its prescribed restoring acceleration. Thus four quarter-periods
give the stated period. For \(n=1\), substitution \(u=\sin\theta\)
gives integral \(\pi/2\), recovering the linear result.
\end{proof}
\end{annexitem}

\begin{annexitem}{annex:p15.approximation.rectangle}{thm:p15-approximation-rectangle}{Height, width and omitted tails}
\begin{proof}
The integral triangle inequality yields
\(\left|\int_I(f-c)\right|\leq\int_I|f-c|\leq dL\).
The same inequality on an omitted region bounds its missing contribution
by the integral of \(|f|\). In the centered Lipschitz case,
\(|f(t)-f(t_0)|\leq K|t-t_0|\), so the error is at most
\(K\int_{-L/2}^{L/2}|s|\,ds=KL^2/4\). These controls, not visual
resemblance, supply the error guarantee.
\end{proof}
\end{annexitem}

\begin{annexitem}{annex:p15.approximation.fall}{thm:p15-approximation-fall}{Quadratic drag and the terminal-speed error}
\begin{proof}
Put \(\tau=gt/V\). Differentiating \(V\tanh\tau\) gives
\(g\operatorname{sech}^2\tau=g(1-\tanh^2\tau)\), verifying the
equation and initial value. Two solutions \(v,w\) with the same initial
value have difference satisfying
\((v-w)'=-g(v+w)(v-w)/V^2\). An integrating factor makes this difference
zero on every common interval, proving uniqueness. Integration gives the
distance because \((\log\cosh\tau)'=\tanh\tau\).

Using \(\log\cosh\tau=\tau+\log(1+e^{-2\tau})-\log2\),
\[
 Vt-s(t)=\frac{V^2}{g}
             [\log2-\log(1+e^{-2\tau})].
\]
For \(\tau\geq0\), the bracket lies in \([0,\log2]\).
Also \(1-\tanh\tau=2/(e^{2\tau}+1)\leq2e^{-2\tau}\).
Finally \(\cosh\tau=e^{gs/V^2}\) and
\(\tanh^2\tau=1-\operatorname{sech}^2\tau\), with nonnegative
velocity, give the distance-based formula. Finite \(\tau\) has
\(\tanh\tau<1\), excluding exact arrival at terminal speed in finite time.
\end{proof}
\end{annexitem}

\begin{annexitem}{annex:p15.approximation.easy-cases}{thm:p15-approximation-easy-cases}{Why endpoint agreement is not a proof}
\begin{proof}
The added polynomial vanishes at every checked point and is positive at
all other points when \(c>0\). Thus the interval admits distinct
candidates with identical checked values. Substitution gives
\(g_c(0)=1\), while its ratio to \(1/x\) is
\[
 \frac{x}{1+x}\left(1+\frac{cx}{(1+x)^2}\right)\longrightarrow1.
\]
At one its value is \((1+c/4)/2\). The free constant makes that value
arbitrarily large while preserving both endpoint conditions. All these
functions are smooth on the domain, so smoothness alone is insufficient.
\end{proof}
\end{annexitem}

\begin{annexitem}{annex:p15.approximation.counting-cases}{thm:p15-approximation-counting-cases}{Inclusive indexing and birthday products}
\begin{proof}
The bijection \(j\mapsto m_1+j\) from
\(\{0,\ldots,m_2-m_1\}\) proves the inclusive count.
For birthdays, expose assignments in order. Conditional on the first
\(j\) being distinct, exactly \(d-j\) choices for the next preserve
distinctness. Independence and uniformity give conditional probability
\((d-j)/d\). Multiplication for \(j=0,\ldots,n-1\) gives the no-match
probability; subtract from one. No assignments give empty product one
and no match. More than \(d\) assignments cannot be distinct by the
pigeonhole principle. Other birthday distributions need different
conditional probabilities.
\end{proof}
\end{annexitem}

\begin{annexitem}{annex:p15.approximation.mechanical-cases}{thm:p15-approximation-mechanical-cases}{Four model-specific determinations}
\begin{proof}
Ellipse slicing gives
\(2b\int_{-a}^a\sqrt{1-x^2/a^2}\,dx\). Substitution
\(x=a\sin\theta\), \(-\pi/2\leq\theta\leq\pi/2\), changes this to
\(2ab\int_{-\pi/2}^{\pi/2}\cos^2\theta\,d\theta=\pi ab\).

The Atwood string constraint gives opposite accelerations of equal
magnitude, and the ideal string/pulley gives common tension. Newton's
equations are \(m_1a=m_1g-T\) and \(-m_2a=m_2g-T\).
Subtracting eliminates tension and gives
\((m_1+m_2)a=(m_1-m_2)g\). Substitute into either equation for tension.
The denominator is positive. Equal masses and either zero-mass limit
reproduce the easy cases, while the equations determine the whole function.

For the projectile, integrating constant acceleration gives
\(X(t)=v\cos\theta\,t\), \(Y(t)=v\sin\theta\,t-gt^2/2\).
For \(v\sin\theta>0\), the positive return time is
\(2v\sin\theta/g\); substituting it gives the range. Horizontal
release from the receiving level has zero flight interval in this model,
and vertical launch has zero horizontal speed. The expression includes
both endpoints and zero launch speed.

For rolling, let \(f\) be static friction directed up the plane.
Translation gives \(ma=mg\sin\theta-f\); rotation and no slip give
\(I\alpha=fr\), \(a=r\alpha\). Hence \(f=Ia/r^2\), and substitution
yields \((m+I/r^2)a=mg\sin\theta\). The required friction must be
available; otherwise no slip fails. Substituting \(I=mr^2\) or
\(I=mr^2/2\) gives the ring and disc accelerations.
\end{proof}
\end{annexitem}

\begin{annexitem}{annex:p15.approximation.dominance}{thm:p15-approximation-dominance}{Positive addition and signed cancellation}
\begin{proof}
For positive summands, \(A+B=M+m=M(1+r)\), \(0<r\leq1\).
Omitting \(m\) has error \(m/(M+m)=r/(1+r)\); the size bounds follow
from the same equation. For signed summands,
\(|A+B|\geq|A|-|B|\geq(1-r)|A|>0\). Dividing the omitted magnitude
\(|B|\leq r|A|\) by that bound gives \(r/(1-r)\). The loss of control
as \(r\to1\) records possible cancellation.
\end{proof}
\end{annexitem}

\begin{annexitem}{annex:p15.approximation.competing-energies}{thm:p15-approximation-competing-energies}{A scale balance and the stationary point}
\begin{proof}
The derivative is \(E'(r)=(Br-2A)/r^3\), negative before \(2A/B\)
and positive afterward. Also \(E\to+\infty\) at zero and
\(E\to0^-\) at infinity, proving the unique global minimum.
Substitution gives \(E_0=-B^2/(4A)\). Equating \(A/r^2\) and \(B/r\)
instead gives \(r=A/B\).

Put \(\epsilon=|E_0|-\Delta\), so
\(0<\epsilon<B^2/(4A)\). Turning points solve
\(\epsilon r^2-Br+A=0\). Its two roots are positive and have sum
\(B/\epsilon\); their midpoint is
\(B/(2\epsilon)=r_0/(1-\Delta/|E_0|)>r_0\).
This determines an extrema midpoint. Time averages weight by time spent
at each position, while thermal expectations require a distribution.
Neither weighting follows from extrema alone. For instance,
\(r(t)=1+2\sin^4t\) has extrema one and three, with midpoint two,
but time mean \(1+2(3/8)=7/4\). The integral follows by applying
\(\sin^2t=(1-\cos2t)/2\) twice over a period. This illustrative path
is not claimed to solve the preceding potential model; it demonstrates
the distinction between an extrema midpoint and an unspecified mean.
\end{proof}
\end{annexitem}

\begin{annexitem}{annex:p15.approximation.water-regimes}{thm:p15-approximation-water-regimes}{The dispersion relation and all corner bounds}
\begin{proof}
For a sinusoidal traveling mode the vertical factor solves
\(W''-k^2W=0\). Its general form is a linear combination of
\(\cosh(k(z+h))\) and \(\sinh(k(z+h))\); the bottom derivative
condition sets the latter coefficient to zero. Thus, choosing a phase,
take a nonzero mode
\[
 \phi=A\cosh(k(z+h))\sin(kx-\omega t),\qquad
 \eta=B\cos(kx-\omega t).
\]
The Laplace equation and bottom condition hold. The surface conditions are
\[
 \omega B=Ak\sinh(kh),\qquad
 A\omega\cosh(kh)=(g+\gamma k^2/\rho)B.
\]
Eliminating the nonzero amplitudes gives
\(\omega^2=(gk+\gamma k^3/\rho)\tanh(kh)\).
The squared phase speed \((\omega/k)^2\) is the stated expression.

For \(t\geq0\), \(0\leq\tanh t\leq t\): its value starts at zero
and its derivative lies in \([0,1]\). Therefore
\[
 0\leq t-\tanh t=\int_0^t\tanh^2s\,ds
                \leq\int_0^ts^2\,ds=t^3/3.
\]
Divide by \(t>0\). The identity
\(1-\tanh t=2/(e^{2t}+1)\) proves the deep-water bound.
The gravity term is \(g/k=g\ell\), while the capillary term is
\(\gamma k/\rho=\gamma/(\rho\ell)\). Multiplication by
\(kh=h/\ell\) gives the two shallow corners; multiplication by one
gives the two deep corners. For \(\gamma>0\), their relative strength
is measured by \(\rho g/(\gamma k^2)\). When \(\gamma=0\), only the
gravity column applies.

Write the exact sum of restoring terms as \(A(1+d)\),
\(0\leq d\leq\delta\), with retained term \(A>0\). Write the exact
depth factor divided by its chosen approximation as \(1-e\),
\(0\leq e\leq\epsilon\), using either preceding bound. The squared
speed ratio is \((1+d)(1-e)\), whose difference from one is bounded by
\(|d-e-de|\leq\delta+\epsilon+\delta\epsilon\).
Keeping both restoring terms and the exact hyperbolic tangent supplies
the intermediate regimes directly from the model equations.
\end{proof}
\end{annexitem}

\begin{annexitem}{annex:p15.approximation.rational}{thm:p15-approximation-rational}{A remainder before a threshold decision}
\begin{proof}
Since \((1+x)(1-x)=1-x^2\), division by the positive denominator gives
the identity and nonnegative remainder \(x^2/(1-x)\). When \(x\leq1/2\),
the denominator is at least one half, proving the upper bound.
The interval certificate then supplies the threshold implications.
At \(x=9/10\), the approximation is \(19/10<5\) and the exact value
is \(10>5\). At \(x=1/10\), \(11/10\) is below \(221/200\) by
\(1/200\), but
\[
 \frac{10}{9}-\frac{221}{200}=\frac{11}{1800}>0.
\]
A smaller approximation error can still exceed a smaller decision margin.
\end{proof}
\end{annexitem}

\begin{annexitem}{annex:p15.approximation.nonuniform}{thm:p15-approximation-nonuniform}{A moving input reaches the transition}
\begin{proof}
The error is exactly \(\varepsilon/(t+\varepsilon)\). At every fixed
\(t>0\), it tends to zero as \(\varepsilon\to0^+\). For each fixed
\(\varepsilon>0\), it is less than one but tends to one as \(t\to0^+\),
so the supremum on \((0,1]\) is one. At the permitted input
\(t=\varepsilon\), the function equals one half. Finally the error
decreases with \(t\), giving the bound for \(t\geq t_0\).
\end{proof}
\end{annexitem}
